\section{Data assimilation -- Smoothing with ES-MDA}
\label{sec:da}

This section describes the data-assimilation strategy implemented in
the rollout script \rolloutscript{}.
We first introduce \ESMDA{} for a single window with a time-varying
parameter vector (Section~\ref{subsec:da-single-window}), then
describe how the prior parameter ensemble is generated for cold and
warm starts (Section~\ref{subsec:da-prior}), and finally explain
how the rollout chains windows together
(Section~\ref{subsec:da-multiwindow}).

\subsection{Single-window ES-MDA for time-varying parameters}
\label{subsec:da-single-window}

Within a single window, the parameter trajectory is represented as a
discrete vector
\begin{equation}
  \paramTraj_j = \bigl(\param_j(t_0), \param_j(t_1), \ldots, \param_j(t_{\Nt-1})\bigr)
  \in \R^{n_p \cdot \Nt}, \qquad j = 1, \ldots, \Ne,
  \label{eq:flat-param}
\end{equation}
where each time point of each physical parameter is treated as an
independent scalar variable.  Stacking time points contiguously
within a parameter -- $(\theta^{(1)}_0, \theta^{(1)}_1, \ldots,
\theta^{(1)}_{\Nt-1}, \theta^{(2)}_0, \ldots)$ -- ensures that
different physical quantities are never mixed during the Kalman
update.  This flatten/unflatten pair is implemented in the
\tvESMDAclass{} class.

\paragraph{Tempering.}
\ESMDA{}~\cite{emerick2013esmda} approximates a single Bayesian
update by a sequence of $\Nsteps$ smaller analyses with inflated
observation error covariance.  The inflation factors
$\{\alphaInfl_i\}_{i=1}^{\Nsteps}$ satisfy
\begin{equation}
  \sum_{i=1}^{\Nsteps} \frac{1}{\alphaInfl_i} = 1,
  \label{eq:esmda-tempering}
\end{equation}
which guarantees that the product of the $\Nsteps$ tempered
Gaussian likelihoods, each with covariance $\alphaInfl_i \CD$,
recovers the original observation likelihood with covariance $\CD$.
The implementation uses the uniform schedule $\alphaInfl_i =
\Nsteps$, so each step contributes equal Bayesian weight.  At each
step the analysis is performed with the standard ensemble Kalman
update below.

\paragraph{Kalman update.}
Let $\paramTraj^{(i)}$ denote the $(n_p \Nt) \times \Ne$ matrix
whose columns are the flattened parameter trajectories of the
ensemble members at the start of step $i$, and $\mat{D}^{(i)}$ the
corresponding $\Nd \times \Ne$ matrix of predicted observations,
$\mat{D}^{(i)}_j = \obsop\bigl(\fwd(\paramTraj^{(i)}_j)\bigr)$.
Define ensemble means and anomalies
\begin{equation}
  \bar{\paramTraj}^{(i)}
    = \frac{1}{\Ne}\sum_{j=1}^{\Ne} \paramTraj^{(i)}_j,
  \quad
  \anomdelta\paramTraj^{(i)} = \paramTraj^{(i)} - \bar{\paramTraj}^{(i)} \mathbf{1}\transp,
  \quad
  \anomdelta\mat{D}^{(i)} = \mat{D}^{(i)} - \bar{\mat{D}}^{(i)} \mathbf{1}\transp,
  \label{eq:anomalies}
\end{equation}
with $\mathbf{1}$ the $\Ne$-vector of ones.  The empirical
cross-covariance between flattened parameters and predicted
observations and the empirical predicted-observation covariance are
\begin{equation}
  \CMD^{(i)} = \frac{1}{\Ne - 1}\,
               \anomdelta\paramTraj^{(i)}\, \bigl(\anomdelta\mat{D}^{(i)}\bigr)\transp,
  \qquad
  \CDD^{(i)} = \frac{1}{\Ne - 1}\,
               \anomdelta\mat{D}^{(i)}\, \bigl(\anomdelta\mat{D}^{(i)}\bigr)\transp.
  \label{eq:cov}
\end{equation}
Each ensemble member is then updated through
\begin{equation}
  \paramTraj^{(i+1)}_j
  = \paramTraj^{(i)}_j
  + \CMD^{(i)} \bigl(\CDD^{(i)} + \alphaInfl_i \CD\bigr)\inv
    \bigl(\obs + \sqrt{\alphaInfl_i}\, \CD^{1/2} \vect{z}_j - \mat{D}^{(i)}_j\bigr),
  \quad \vect{z}_j \sim \mathcal{N}(\zerovec, \identity).
  \label{eq:esmda-update}
\end{equation}
Crucially, the ensemble forecast for step $i{+}1$ re-runs the
forward model from the same window-initial state $\state\win{w}(0)$
with the updated parameters.  Without this reset, iteration $i{+}1$
would forecast from iteration $i$'s output, which would bypass the
spin-up after the first iteration and entangle the iterates'
physical states.

\paragraph{Pinning the initial time point.}
The Kalman update may optionally be restricted to time indices
$t_1, \ldots, t_{\Nt-1}$ within the window, leaving $\param_j(t_0)$
untouched.  This is controlled by the \code{pin\_initial\_time\_point}
flag in \tvESMDAclass{}: when set, the entries
corresponding to $t_0$ are dropped from the flattened augmented
state in~\eqref{eq:flat-param} and reinserted unchanged on
unflattening.  Pinning is essential for the rollout
(Section~\ref{subsec:da-multiwindow}) because the boundary value
$\param_j(t_0)$ at window $w > 0$ is set by the previous window's
posterior; any further adjustment by \ESMDA{} would break
continuity at the window seam.

\subsection{Prior parameter generation}
\label{subsec:da-prior}

The prior parameter ensemble is the entry point of every \ESMDA{}
window.  In \pyurbanair{} the prior is provided by a
\paramTSclass{} object, which exposes two methods:
\begin{itemize}
  \item \code{sample\_prior(time\_coords, rng\_key)} -- a cold-start
        draw used for window $0$ and, in identical-twin experiments,
        also for the truth (with $\Ne = 1$);
  \item \code{extrapolate(posterior, prediction\_times, rng\_key)} --
        a warm-start propagation of a window's posterior into the
        prior of the next window.
\end{itemize}
Four implementations of \paramTSclass{} are registered under
\code{src/\allowbreak pyurbanair/\allowbreak parameter\_time\_series/}.
All four share an
external prior specification $(\xext, \Sext)$ per parameter -- in
the notation of Evensen~\cite{evensen2024}, the climatological mean
and standard deviation of each parameter -- together with a
correlation length $\ell$ that controls the smoothness of the
trajectory in time.

\paragraph{Gaussian process with linear-trend extrapolation
(\code{gp\_linear\_trend}).}
The cold-start prior is a smooth ensemble drawn from an RBF Gaussian
process~\cite{rasmussen2006gp} with marginal $\mathcal{N}(\xext,
\Sext^2)$ and correlation length $\ell$.  Extrapolation fits a
per-member linear trend to the posterior trajectory by ordinary
least squares and adds a zero-mean GP residual with heteroscedastic
noise estimated from the cross-ensemble spread.  An optional
exponential damping prevents the trend from diverging far into the
next window.

\paragraph{Stationary AR(1) (\code{ar1}).}
The cold-start prior is a stationary first-order auto-regressive
process with $\phi = e^{-\Delta t / \ell}$ around $\xext$.
Extrapolation fits a per-member AR(1) coefficient by ordinary least
squares and rolls the trajectory forward deterministically from each
member's last training value.

\paragraph{Ornstein--Uhlenbeck (\code{ornstein\_uhlenbeck}).}
The cold-start prior is a stationary Ornstein--Uhlenbeck
process~\cite{uhlenbeck1930ou}.  Extrapolation fits an
\OU{}-with-drift to the posterior, advances it stochastically with
an Euler--Maruyama step, and pools the diffusion coefficient across
the ensemble to keep the spread calibrated.

\paragraph{Critically-damped AR(2) relaxation (\code{ar2\_relaxation}).}
The default in our experiments and the model used in
Evensen~\cite{evensen2024}.  The latent process satisfies the
critically-damped second-order SDE
\begin{equation}
  \frac{\mathrm{d} z}{\mathrm{d} t} = w,
  \qquad
  \frac{\mathrm{d} w}{\mathrm{d} t} = -2\lambda\, w - \lambda^2 z + \eta(t),
  \qquad
  \lambda = \frac{\sqrt{3}}{\ell},
  \label{eq:ar2}
\end{equation}
where $\eta$ is a white-noise term scaled so that $z$ has zero mean
and unit variance in stationarity.  The implementation
integrates~\eqref{eq:ar2} with the closed-form transition matrix
$\mat{F} = \exp(\mat{A}\,\Delta t)$ and the exact one-step process
noise $\mat{Q} = \mat{P}_{\mathrm{stat}} - \mat{F}\,
\mat{P}_{\mathrm{stat}}\, \mat{F}\transp$ with $\mat{P}_{\mathrm{stat}}
= \diag(1, \lambda^2)$, factored via Cholesky.  Each grid step
preserves the stationary covariance exactly, without substepping.

\medskip
\noindent\emph{Cold start.}\;
For window $0$ the prior follows
Evensen~\cite{evensen2024}:
\begin{equation}
  \param_j(t) = \xext + \Sext\, z_j(t),
  \label{eq:ar2-cold}
\end{equation}
with $z_j$ drawn from the stationary distribution of~\eqref{eq:ar2}.

\medskip
\noindent\emph{Warm start.}\;
For window $w > 0$ let $\mu_{\mathrm{end}} = \Ne^{-1}
\sum_j \param_j^{\post,(w-1)}(t_{\Nt-1})$ be the ensemble-mean
end-of-window posterior of the previous window, and $z_j(t)$ a fresh
trajectory of~\eqref{eq:ar2} on the new window's grid.  The next
window's prior is then constructed by exponential relaxation toward
the external prior:
\begin{equation}
  \param_j(t) = \alpha(t)\, \mu_{\mathrm{end}}
              + \bigl(1 - \alpha(t)\bigr)\,\bigl(\xext + \Sext\, z_j(t)\bigr),
  \quad
  \alpha(t) = e^{-(t - t_0)/\ell}.
  \label{eq:ar2-relax}
\end{equation}
At $t = t_0$ the relaxation coefficient $\alpha = 1$ and every
member starts at the previous-window ensemble mean
$\mu_{\mathrm{end}}$, so the per-member spread at the seam is zero
and the seam is determined by the previous posterior's mean.  Far
into the window $\alpha(t) \to 0$ and the prior relaxes back to
$\xext$ with full external spread $\Sext$.  The AR(2) state $(z, w)$
is initialized per member from the normalized end-of-window
posterior so that the latent process is itself $C^1$-continuous
across the window seam, even though only the ensemble mean appears
explicitly in~\eqref{eq:ar2-relax}.

\paragraph{Avoiding the inverse crime.}
In identical-twin experiments, using the same generative process for
the truth and for the prior would underestimate the difficulty of the
inverse problem.  The truth trajectory is therefore drawn from a
separate \paramTSclass{} instance, typically with a
different correlation length, that samples once over the full
horizon $[0, \Nw\windowLen]$ and slices the result into per-window
datasets, sharing boundary points so that the truth is continuous
across windows.

\subsection{Multi-window rollout}
\label{subsec:da-multiwindow}

The rollout chains the single-window \ESMDA{} of
Section~\ref{subsec:da-single-window} with the prior generators of
Section~\ref{subsec:da-prior} into a sequence covering the full
horizon.  Algorithm~\ref{alg:rollout} summarizes the loop, which
mirrors the implementation in
\rolloutscript{}.

\begin{algorithm}[t]
\caption{Rolling-window \ESMDA{} with time-varying parameters.}
\label{alg:rollout}
\begin{algorithmic}[1]
\REQUIRE Number of windows $\Nw$, window length $\windowLen$, time
points per window $\Nt$, ensemble size $\Ne$, external prior
$(\xext, \Sext)$, time-series model \code{ts\_model}, truth model
\code{truth\_ts\_model}, forward models $\fwd_{\mathrm{truth}}$ and
$\fwd_{\mathrm{assim}}$, \ESMDA{} configuration $(\Nsteps, \CD)$.
\STATE \emph{(Truth pre-generation.)}\;
       Draw a single trajectory from \code{truth\_ts\_model} on a
       shared-boundary grid covering $[0, \Nw\windowLen]$ and slice
       into $\{\paramTraj^{\star\winSub{w}}\}_{w=0}^{\Nw-1}$, each
       on local time $[0, \windowLen]$.
\STATE \emph{(Initial prior.)}\;
       $\paramTraj_{\mathrm{ens}} \leftarrow$
       \code{ts\_model.sample\_prior}$\bigl([0, \windowLen],\,
       \cdot\bigr)$.
\STATE $\state_\star \leftarrow \mathtt{None}$;\quad
       $\state_{\mathrm{end}} \leftarrow \mathtt{None}$.
\FOR{$w = 0, 1, \ldots, \Nw - 1$}
  \STATE Record the prior used for window $w$:
         $\paramTraj^{\prior,(w)} \leftarrow \paramTraj_{\mathrm{ens}}$.
  \STATE Run the truth solver:
         $\state_\star \leftarrow \fwd_{\mathrm{truth}}\bigl(
         \paramTraj^{\star\winSub{w}};\, \state_\star\bigr)$.
  \STATE Generate noisy observations
         $\obs\win{w} = \obsop_{\mathrm{truth}}(\state_\star)
         + \CD^{1/2}\,\vect{z}$,
         $\vect{z} \sim \mathcal{N}(\zerovec, \identity)$.
  \STATE \textbf{if} $w > 0$ \textbf{then} enable
         \code{pin\_initial\_time\_point} so the seam value
         established by \code{extrapolate} is preserved.
  \STATE Choose the warm-start state for the assimilation forecast:
         $\state\win{w}(0) =
         \state_{\mathrm{end}}\big|_{t=\windowLen}$ when $w > 0$,
         else $\mathtt{None}$ (cold start with spin-up).
  \STATE Run \ESMDA{} with $\Nsteps$ tempering steps on the
         flattened parameter vector, using
         $\fwd_{\mathrm{assim}}$ for the forecasts; obtain the
         parameter history $\paramTraj^{(0)}, \ldots,
         \paramTraj^{(\Nsteps)}$ and the final ensemble forecast
         state $\state_{\mathrm{end}}\win{w}$.
  \STATE Record posterior $\paramTraj^{\post,(w)} \leftarrow
         \paramTraj^{(\Nsteps)}$.
  \IF{$w < \Nw - 1$}
    \STATE Build next window's prior on $[\windowLen, 2\windowLen]$
           and relabel to local $[0, \windowLen]$:\;
           $\paramTraj_{\mathrm{ens}} \leftarrow$
           \code{ts\_model.extrapolate}$\bigl(
           \paramTraj^{\post,(w)};\,
           [\windowLen, 2\windowLen],\, \cdot\bigr)$.
  \ENDIF
\ENDFOR
\ENSURE Per-window priors $\{\paramTraj^{\prior,(w)}\}$, posteriors
$\{\paramTraj^{\post,(w)}\}$, truth states, and assimilation states;
concatenated on absolute time at save time.
\end{algorithmic}
\end{algorithm}

\paragraph{Local-time invariance.}
Each window's simulation runs with its own clock $[0, \windowLen]$,
independently of $w$.  The extrapolation in step~12 is performed on
$[\windowLen, 2\windowLen]$ -- the natural prediction interval given
the previous posterior's local time grid -- and then relabeled to
$[0, \windowLen]$ for the next window's forward model.  This keeps
the smoother's \code{time\_coords} constant across windows and
spares the forward model from recomputing time stamps; it works for
all four prior models because the AR(2) relaxation is
translation-invariant in local time, and the GP- and AR(1)-based
methods only depend on relative offsets within their training and
prediction grids.

\paragraph{Continuity at boundaries.}
Continuity across windows is enforced by two complementary
mechanisms.  First, the prior generator -- in the case of the AR(2)
relaxation, equation~\eqref{eq:ar2-relax} -- guarantees a well-defined
seam value $\param_j\win{w+1}(t_0)$ inherited from window $w$.
Second, pinning the initial time point in \ESMDA{} from window $1$
onward ensures that the Kalman update preserves that seam value.
Without pinning, the smoother would treat $\param_j(t_0)$ as a free
parameter and the trajectory would jump at every window boundary.

\paragraph{Storage.}
The rollout records, for every window, the prior parameter ensemble,
the posterior parameter ensemble, the truth state, and the final
ensemble forecast state.  These are concatenated into global NetCDF
archives after shifting from local to absolute time, and they form
the basis for the diagnostics shown in
Section~\ref{sec:results}.
